% !TEX root = ../main.tex

\chapter{引言}

研究生学位论文质量反映了研究生培养质量，是授予相应学位的主要依据。
学位论文要能代表研究生研究工作的水平，也是申请和授予相应学位的主要依据。
提高学位论文的质量和水平是研究生教育最重要的任务，需要学生和导师的共同努力。
通常学位论文只能有一个主题，该主题应针对某学科领域中的一个具体问题展开深入、
系统的研究，并得出有价值的研究结论。


\section{内容要素}

\subsection{题名页}

题名页主要包括以下内容：

A) 分类号：需同时标注中图分类号与UDC分类号两种。
选择分类号时，应以论文内容主题所涉及的科学领域为依据，
可同时参考自身专业。一般情况下选至二级或三级类目即可，
也可选至更精确的类目号。
若不确定具体分类，可咨询学校图书馆。

a) 中图分类号：采用《中国图书馆分类法》最新版本标注，
可参考：\url{https://www.clcindex.com/category}。
字符总数对应类目的级数，例如：论文中涉及火星，专业为地球物理学，
二级分类号为P1（天文学），三级分类号为P18（太阳系），
或者更详细至P185.3（火星）。

b) UDC分类号：采用《国际十进分类法》（Universal Decimal
Classification） 最新版本标注。
可参考：\url{https://udcsummary.info/php/index.php?lang=chi}。
数字总数对应类目的级数，例如：论文中涉及行星地质学，专业为地球物理学，
二级分类号为55（地球科学、地质学），三级分类号为550（地质学附属学科），
四级分类号为550.3（地球物理学）。

B) 学校代码：中国科学技术大学的学校代码为10358。

C) 密级：按GB/T 7156的规定进行标注。仅限涉密学位论文在题名页须做出国家秘密标志，
其他论文可不标注。

D) 题名：题名以简明的词语恰当、准确地反映论文最重要的特定内容，一般不宜超过25字，
应中英文对照。

E) 责任者：责任者包括学生姓名、学号，指导教师姓名、专业技术职务等。
校外导师应一并列出。由论文指导委员会指导论文写作的，应列出全部指导教师。

F) 学科专业：参照国务院教育行政部门颁布的最新《研究生教育学科专业目录》
标注硕士和博士学位论文的学科和专业。

G) 研究方向：指本学科专业范畴下的细分研究领域。

H) 论文提交日期：指论文上交日期（不得晚于学位授予日期）。



\subsection{中国科学技术大学学位论文原创性和使用授权声明}

本部分内容使用统一的模板，具体内容见格式范例。
授权使用声明须仔细阅读后慎重选择公开或控阅（控阅年份必填）。
提交学位论文纸质版和电子版时，须有学生和导师签名（手签或电子签，
但不能出现签章平台的二维码），同时须有签字日期（不得晚于学位授予日期）。


\subsection{研究生学位论文（实践成果报告）使用（或未使用）人工智能工具声明}

本部分内容使用统一的模板，具体内容见格式范例。
提交学位论文纸质版和电子版时须有作者签名（手签或电子签，
但不能出现签章平台的二维码）、日期（不得晚于学位授予日期）。
作者须根据实际情况如实声明：若在研究或撰写过程中使用了人工智能工具，
应保留“使用人工智能工具声明”页，并详细披露所使用的生成式人工智能工具名称及具体用途。
若未使用人工智能工具，请选择保留“未使用人工智能工具声明”页
（注：未选用的另一份声明页需在定稿中删除）。


\subsection{插图和附表清单}

论文中插图和附表较多时，应分别列出“插图清单”和“附表清单”。
插图清单在前，应列出图序、图题和页码。
附表清单在后，应列出表序、表题和页码。
插图较多而附表较少、或者插图较少而附表较多、或者二者均较少时，
可将插图和附表合在一起列出“插图和附表清单”，插图在前、附表在后。


\section{撰写要求}

\section{论文的基本要求}

硕士学位论文正文要求一般不少于3万字，博士学位论文正文要求一般不少于5万字。
留学生学位论文字数参照以上篇幅综合考虑。

写作时要注意论文具有正确的政治、思想导向，一定的知识性、科学性以及原创性，
论文的内容、体例与文字等符合现行规范。

其中，工程博士学位论文中必须至少有专门一章或一节介绍学位论文研究内容的工程
应用现状和前景。
